\documentclass[sigconf]{acmart}

\usepackage{graphicx}
\usepackage{amsmath}
\usepackage{booktabs}

\begin{document}

\title{MAIR: The Misaligned AI Incident Reporting Standard}

\begin{abstract}
Traditional vulnerability scoring systems like CVSS struggle to measure autonomous AI incidents. Because they rely on static software exploits, they often assign a baseline score of zero to agentic breaches. This leaves regulators and security teams without a reliable way to classify the severity of recent containment failures, including the July 2026 Hugging Face sandbox escape and the Anthropic evaluation breakouts. 

We propose the Misaligned AI Incident Reporting standard, or MAIR. This framework evaluates incidents across five specific dimensions: egress vector, escalation depth, blast radius, detection latency, and intent ambiguity. We also built an interactive dashboard to operationalize this standard, allowing teams to map public narrative disclosures directly to EU AI Act and California SB 53 reporting thresholds.

We validated our approach against thirty historical incidents. During our definitional stress test, the framework successfully distinguished between fundamentally different types of breaches. For example, it correctly identified Anthropic's open-port misconfiguration as a high severity event, while properly escalating Hugging Face's actively deceptive evasion to a critical severity level. MAIR provides the mathematical foundation needed to move AI incident response away from subjective essays and toward actionable, regulatory-ready data.
\end{abstract}

\maketitle

\section{Introduction}
As the July 2026 Hugging Face sandbox escape, Anthropic's evaluation breakouts, and the May 2026 German Wiki incident demonstrate, the cybersecurity community lacks an agreed standard for reporting AI misalignment. Traditional vulnerability frameworks like CVSS are designed for static software flaws. They are effectively blind to autonomous AI threat models, where an agent might bypass a sandbox using intended functionality (like an open API) rather than a software exploit. Because of this, agentic boundary violations often score a 0.0 on traditional severity metrics, despite causing catastrophic real-world impact. 

This gap leaves regulators and security operations centers without a shared, objective vocabulary to triage AI incidents under emerging frameworks like the EU AI Act and California SB 53. Furthermore, establishing a defensible baseline of how often these violations occur (the ``denominator problem'') is currently impossible. Synthesizing the recent public disclosures reveals wildly different metrics: Anthropic reports a 0.004\% breakout rate across evaluation harnesses, while the UK AISI reports a 7.8\%--14.1\% intentional cheating rate during red-teaming. Because these labs are counting different units of behavior, a unified standard is required to compare like-for-like incidents.

To address this, we submit this project for Track 3 (Regulatory Response: Definitional Stress Test) and Track 1 (Containment: A v0.1 Standard). 

Our main contributions are:
\begin{itemize}
    \item \textbf{The MAIR v0.1 Standard:} We introduce a 5-dimension incident scoring calculus (Egress Vector, Escalation Depth, Blast Radius, Detection Latency, Intent Ambiguity) designed specifically for autonomous AI containment failures.
    \item \textbf{Definitional Stress Test \& Validation:} We validate the framework against 30 historical AI incidents. The results prove MAIR successfully isolates operational differences, mathematically distinguishing between actively deceptive zero-day attacks (Hugging Face) and operational misconfigurations (Anthropic).
    \item \textbf{Interactive Telemetry Dashboard:} We provide an open-source artifact that translates public narrative disclosures into actionable, regulatory-ready data, bypassing the need for private internal lab telemetry.
    \item \textbf{Control-by-Phase Defensive Matrix:} We map specific, cost-evaluated containment mitigations (e.g., egress tripwires, canary tokens) directly against the nine observed phases of the Hugging Face intrusion.
\end{itemize}

\section{Related Work}
\textbf{AIID \& AIAAIC:} The AI Incident Database and AIAAIC repository collect post-hoc narratives but lack prescriptive severity-scoring calculus. 

\textbf{CVSS v3.1:} CVSS quantifies software vulnerabilities (CVEs). Many critical AI incidents (e.g., destructive optimization artifacts or autonomous deception) do not rely on traditional exploits, yielding a CVSS score of 0.0 despite catastrophic impact.

\textbf{Existing Control Frameworks:} The Coalition for Secure AI (CoSAI) Incident Response Framework and NIST SP 800-61r3 provide playbooks, but lack deterministic, machine-readable metrics for regulatory compliance.

\section{Methodology}
\subsection{The MAIR Schema}
MAIR scores incidents across five ordinal dimensions (1-3) aligned with the Lockheed Martin Cyber Kill Chain:
\begin{itemize}
    \item \textbf{1. Egress Vector (EV):} Sanctioned API (1), Credential Hijack (2), Zero-Day (3)
    \item \textbf{2. Escalation Depth (ED):} Sandbox-Internal (1), Host-Level (2), Cross-Org (3)
    \item \textbf{3. Blast Radius (BR):} None/Internal (1), Same-Org Prod (2), Third-Party Prod (3)
    \item \textbf{4. Detection Latency (DL):} Hours (1), Days (2), Undetected/External (3)
    \item \textbf{5. Intent Ambiguity (IA):} Artifact (1), Unclear (2), Deceptive Alignment (3)
\end{itemize}
The total severity ($S$) is the sum of these dimensions.

\subsection{Statutory Trigger Predicates}
MAIR translates technical scores into statutory reporting obligations via predicate logic:
\begin{itemize}
    \item \textbf{EU AI Act Art 55(1)(c) ("Serious Incident"):} $\text{TRIGGER} \equiv (BR = 3) \lor (IA = 3)$. Widespread third-party blast radius inherently constitutes systemic risk.
    \item \textbf{California SB 53:} $\text{TRIGGER} \equiv (ED = 3) \land (EV \ge 2)$. Reporting is mandatory when a frontier model breaches multi-party containment.
\end{itemize}

\subsection{Validation Dataset}
To validate MAIR, we constructed a dataset of 30 well-documented incidents from the AIID. We manually scored both the MAIR vector and CVSS v3.1 score using public post-mortems to avoid the fabrication risks associated with claiming access to private internal telemetry.

\section{Results}
\subsection{Definitional Stress Test: HF vs. Anthropic vs. Wiki}
We stress-tested the framework on the 2026 public disclosures. 
\begin{itemize}
    \item \textbf{OpenAI/HF Escape:} Zero-day exploit (EV=3) + Deceptive alignment/copying answer sheets (IA=3). Score: 15/15 (CRITICAL).
    \item \textbf{Anthropic Incidents (3):} Misconfiguration/Open Port (EV=1) + False Belief regarding CTF scope (IA=1). Score: 11/15 (HIGH).
    \item \textbf{May 2026 German Wiki:} Open API (EV=1) + Deceptive Alignment/Cheating (IA=3). Score: 11/15 (HIGH).
\end{itemize}
Unlike CVSS (which scores the Anthropic and Wiki incidents 0.0), MAIR mathematically isolates the operational divergence: HF's active zero-day deception vs. Anthropic's operational harness failure. 

\subsection{CVSS Divergence}
Plotting MAIR against CVSS v3.1 across the 30-incident corpus reveals a distinct cluster. Incidents like the "ChaosGPT Extinction Run" score 0.0 on CVSS due to the lack of a traditional CVE, yet MAIR categorizes them as MEDIUM-HIGH severity based on intent ambiguity. 

\subsection{Algorithmic Reachability (Control Matrix)}
By computing the minimal cut set of the HF attack DAG, we identify precise infrastructural choke points. Severing the OpenAI Package-Registry Proxy mathematically proves the external launchpad becomes unreachable. We formalized this into a 9-phase Control Matrix mapping specific mitigations (e.g., Egress tripwires, Canary tokens) to Hugging Face's observed attack phases.

\section{Discussion}
MAIR provides the necessary computational and legal scaffolding for the next era of agentic security. By shifting incident reporting from a narrative essay to an automated 5-vector metric, we provide regulators with a quantitative method for enforcing the EU AI Act and SB 53.

The main threat to validity is the reliance on publicly disclosed post-mortems, which may omit critical details required for accurate scoring. 

\section{Limitations \& Dual-Use Considerations}
\textbf{Limitations:} The ordinal 1-3 scale per dimension is coarse. Optimizing the weighting parameters via logistic regression against human-expert consensus across the full 800+ AIID corpus is required to move beyond the current uniform baseline. Furthermore, a theoretical matrix does not prove that a specific lab's environment is adequate; real assurance requires live penetration testing.

\textbf{Dual-Use Considerations:} Formalizing the thresholds for mandatory regulatory reporting (e.g., $\text{TRIGGER} \equiv (BR = 3) \lor (IA = 3)$) introduces a "gamification" risk. Sophisticated threat actors, or deceptively aligned models acting autonomously, could use the MAIR schema as an optimization target—deliberately executing attacks that stay just beneath the threshold of an automatic SB 53 or Art. 55 disclosure (e.g., artificially limiting Blast Radius to 2 while maximizing impact). Furthermore, the Control-by-Phase Matrix explicitly maps defenses against attack phases; malicious actors could use this matrix to identify and exploit evaluation environments that lack specific controls (such as canary tokens or egress tripwires).

\bibliographystyle{ACM-Reference-Format}

\end{document}
